% statistiques-probabilites2.tex --- exemples texecole, niveau collège
% !TeX program = lualatex
% Compilation : lualatex --shell-escape (les calculs exécutent Python)
\documentclass[marges=20, taille=11, precision=3]{texecole}
\begin{document}

soit h1 = <bleu marine gras section num>
soit h2 = <bleu italique sous-section>

<h1>Statistiques et probabilités — collège

<h2>Moyenne, médiane, étendue, quartiles

<Calcule>la moyenne de la série 12 ; 15 ; 9 ; 14

<Calcule>la moyenne de la série de valeurs 8 ; 12 ; 15 et d'effectifs 2 ; 5 ; 3

<Calcule>la médiane de la série 9 ; 15 ; 12 ; 14

<Calcule>la médiane de la série 7 ; 3 ; 9 ; 5 ; 11

<Calcule>l'étendue de la série 12 ; 15 ; 9 ; 14

<Calcule>les quartiles de la série 2 ; 4 ; 5 ; 7 ; 8 ; 10 ; 12 ; 15

<h2>Proportionnalité

<Calcule>la quatrième proportionnelle de 3, 5 et 12

<Vérifie>si le tableau est de proportionnalité {
	2 ; 3 ; 5
	6 ; 9 ; 15
}

<Vérifie>si le tableau est de proportionnalité {
	2 ; 3 ; 5
	6 ; 9 ; 14
}

<Complète>le tableau de proportionnalité {
	2 ; 3 ; ?
	6 ; ? ; 15
}

<h2>Pourcentages

<Calcule>30 % de 250

<Applique>une augmentation de 5 % à 240

<Applique>une diminution de 12 % à 60

<Calcule>le taux d'évolution de 250 à 280

<page suivante h2>Échelles et vitesses

<Calcule>l'échelle d'un plan où 2 cm représentent 50 m

<Calcule>la vitesse moyenne pour 150 km en 2 h 30 min

<Calcule>la distance parcourue à 80 km/h pendant 1 h 45 min

<Calcule>la durée du trajet de 210 km à 60 km/h

\end{document}
